% ============================================================================
% Section III: Mathematical Formulation
% ============================================================================

\section{Mathematical Formulation}
\label{sec:formulation}

\subsection{Sets and Indices}

The GTEP problem is defined over the following sets:

\begin{itemize}
  \item $\setS$: scenarios, indexed by $s$, with probability $\pi_s$
  \item $\setT$: stages (planning periods), indexed by $t$
  \item $\setB_t$: blocks (operating periods) in stage $t$, indexed by $b$,
        with duration $\Delta_b$ [h]
  \item $\setN$: buses (electrical nodes), indexed by $n$
  \item $\setG$: generators, indexed by $g$, with bus $n_g$
  \item $\setD$: demands, indexed by $d$, with bus $n_d$
  \item $\setL$: transmission lines, indexed by $\ell$, connecting buses
        $a_\ell$ and $b_\ell$
  \item $\setE$: batteries (energy storage), indexed by $e$
  \item $\setJ$: hydraulic junctions, indexed by $j$
  \item $\setR$: reservoirs, indexed by $r$, with junction $j_r$
  \item $\setW$: waterways, indexed by $w$
  \item $\setH$: turbines, indexed by $h$
  \item $\setZ$: reserve zones, indexed by $z$
\end{itemize}

\begin{figure}[t]
  \centering
  \begin{tikzpicture}[
    scenbox/.style={draw, rounded corners=6pt, fill=blue!5,
      minimum width=7.6cm, inner sep=4pt},
    stagebox/.style={draw, rounded corners=4pt, fill=blue!10,
      minimum width=3.4cm, inner sep=3pt},
    blockbox/.style={draw, rounded corners=2pt, fill=blue!18,
      minimum width=0.55cm, minimum height=0.55cm,
      font=\scriptsize, inner sep=1pt},
    lbl/.style={font=\scriptsize\itshape, text=black!70},
  ]
    % --- Scenario 1 ---
    \node[blockbox] (s1t1b1) at (0,0) {$b_1$};
    \node[blockbox, right=1pt of s1t1b1] (s1t1b2) {$b_2$};
    \node[blockbox, right=1pt of s1t1b2] (s1t1b3) {$b_3$};
    \node[stagebox, fit=(s1t1b1)(s1t1b3), label={[lbl]above:{Stage $t{=}1$, $\delta_1$}}] (s1t1) {};

    \node[blockbox, right=12pt of s1t1b3] (s1t2b1) {$b_1$};
    \node[blockbox, right=1pt of s1t2b1] (s1t2b2) {$b_2$};
    \node[blockbox, right=1pt of s1t2b2] (s1t2b3) {$b_3$};
    \node[blockbox, right=1pt of s1t2b3] (s1t2b4) {$b_4$};
    \node[stagebox, fit=(s1t2b1)(s1t2b4), label={[lbl]above:{Stage $t{=}2$, $\delta_2$}}] (s1t2) {};

    \node[scenbox, fit=(s1t1)(s1t2),
      label={[lbl]above:{Scenario $s{=}1$, $\pi_1$}}] (sc1) {};

    % --- Scenario 2 ---
    \node[blockbox, below=28pt of s1t1b1] (s2t1b1) {$b_1$};
    \node[blockbox, right=1pt of s2t1b1] (s2t1b2) {$b_2$};
    \node[blockbox, right=1pt of s2t1b2] (s2t1b3) {$b_3$};
    \node[stagebox, fit=(s2t1b1)(s2t1b3), label={[lbl]above:{Stage $t{=}1$}}] (s2t1) {};

    \node[blockbox, right=12pt of s2t1b3] (s2t2b1) {$b_1$};
    \node[blockbox, right=1pt of s2t2b1] (s2t2b2) {$b_2$};
    \node[blockbox, right=1pt of s2t2b2] (s2t2b3) {$b_3$};
    \node[blockbox, right=1pt of s2t2b3] (s2t2b4) {$b_4$};
    \node[stagebox, fit=(s2t2b1)(s2t2b4), label={[lbl]above:{Stage $t{=}2$}}] (s2t2) {};

    \node[scenbox, fit=(s2t1)(s2t2),
      label={[lbl]above:{Scenario $s{=}2$, $\pi_2$}}] (sc2) {};

    % --- Duration label ---
    \draw[decorate, decoration={brace, amplitude=3pt, mirror}]
      (s1t1b1.south west) -- (s1t1b1.south east)
      node[midway, below=3pt, font=\tiny] {$\Delta_b$};
  \end{tikzpicture}
  \caption{Time hierarchy: scenarios $\setS$ contain stages $\setT$,
    each with blocks $\setB_t$.  Probabilities $\pi_s$, discount factors
    $\delta_t$, and block durations $\Delta_b$ weight the objective.}
  \label{fig:time_hierarchy}
\end{figure}

\subsection{Decision Variables}

\noindent\textbf{Operational variables} (per scenario $s$, stage $t$, block $b$):
\begin{align}
  p_{g,s,t,b}  &\in [\underline{p}_g, \overline{p}_g]
    & \text{generator output [MW]} \label{eq:pgen} \\
  \ell_{d,s,t,b} &\in [0, \overline{\ell}_d]
    & \text{served demand [MW]} \label{eq:load} \\
  f_{d,s,t,b}  &\geq 0
    & \text{unserved demand [MW]} \label{eq:fail} \\
  \phi_{\ell,s,t,b} &\in [-\overline{\phi}^{ba}_\ell, \overline{\phi}^{ab}_\ell]
    & \text{line flow [MW]} \label{eq:flow} \\
  \theta_{n,s,t,b} &\in \R
    & \text{voltage angle [rad]} \label{eq:theta}
\end{align}

\noindent\textbf{Storage state variables}:
\begin{align}
  e_{e,s,t,b}  &\in [\underline{e}_e, \overline{e}_e]
    & \text{battery SoC [MWh]} \label{eq:soc} \\
  v_{r,s,t,b}  &\in [\underline{v}_r, \overline{v}_r]
    & \text{reservoir volume [hm\textsuperscript{3}]} \label{eq:vol}
\end{align}

\noindent\textbf{Investment variables} (per stage $t$):
\begin{equation}
  m_{g,t} \in [0, \overline{m}_g]
  \quad \text{expansion modules built} \label{eq:expmod}
\end{equation}

\subsection{Objective Function}

Minimize total expected discounted cost:
\begin{equation}
\label{eq:objective}
\min \sum_{s \in \setS} \pi_s \sum_{t \in \setT} \delta_t
\left[
  \sum_{b \in \setB_t} \Delta_b \cdot C^\text{op}_{s,t,b}
  + C^\text{inv}_t
\right]
\end{equation}
where $\delta_t = (1+r)^{-\tau_t/8760}$ is the discount factor with annual
rate $r$ and $\tau_t$ the cumulative hours to stage $t$.

The operational cost for each $(s,t,b)$ is:
\begin{equation}
\label{eq:opcost}
C^\text{op}_{s,t,b} = \sum_{g \in \setG} c_g \, p_{g,s,t,b}
  + \sum_{d \in \setD} c^f_d \, f_{d,s,t,b}
  + \sum_{\ell \in \setL} c^\phi_\ell \, |\phi_{\ell,s,t,b}|
\end{equation}

The investment cost for stage $t$ is:
\begin{equation}
\label{eq:invcost}
C^\text{inv}_t = \sum_{g \in \setG} \kappa_g \, m_{g,t}
  + \sum_{d \in \setD} \kappa_d \, m_{d,t}
  + \sum_{\ell \in \setL} \kappa_\ell \, m_{\ell,t}
  + \sum_{e \in \setE} \kappa_e \, m_{e,t}
\end{equation}
where $c_g$ is the variable generation cost [\$/MWh], $c^f_d$ is the
demand failure penalty [\$/MWh], $c^\phi_\ell$ is the transmission cost
[\$/MWh], and $\kappa_{(\cdot)}$ is the annualized investment cost
[\$/year] for each component.  All objective coefficients are divided by a
scaling factor $\sigma$ (\texttt{scale\_objective}, default 1000) to
improve solver numerical conditioning.  Concretely, each LP cost
coefficient is weighted by $\omega_{s,t,b} = \pi_s \delta_t \Delta_b /
\sigma$, so the reported objective equals $z/\sigma$.

\subsection{Bus Power Balance}

At every bus $n$, for each $(s,t,b)$:
\begin{multline}
\label{eq:balance}
\sum_{\substack{g \in \setG \\ n_g = n}} (1{-}\lambda_g) p_{g,s,t,b}
- \sum_{\substack{d \in \setD \\ n_d = n}}
  (1{+}\lambda_d)(\ell_{d,s,t,b} + f_{d,s,t,b}) \\
+ \sum_{\ell \in \setL} A_{n,\ell} \, \phi_{\ell,s,t,b} = 0
\end{multline}
where $\lambda_g, \lambda_d$ are injection and withdrawal loss factors,
respectively.  The line flow terms use directional contributions rather
than a simple incidence matrix: for each line~$\ell$ from bus~$a_\ell$
to~$b_\ell$, the sending bus sees $-\phi^+_\ell$ (or $+(1-\lambda_\ell)
\phi^-_\ell$ when losses are modeled) and the receiving bus sees
$+(1-\lambda_\ell)\phi^+_\ell$ (or $-\phi^-_\ell$), where $\lambda_\ell$
is the line loss factor.  When $\lambda_\ell=0$, a single bidirectional
variable suffices and the incidence reduces to $A_{n,\ell} \in
\{+1,-1,0\}$.  The dual of~\eqref{eq:balance} is the locational marginal
price (LMP) at bus~$n$.

\subsection{DC Power Flow (Kirchhoff Constraints)}

When \texttt{use\_kirchhoff} is enabled, for each line $\ell$ and
$(s,t,b)$:
\begin{equation}
\label{eq:kirchhoff}
\phi_{\ell,s,t,b} = \frac{V_\ell^2}{X_\ell}
  \left(\theta_{a_\ell,s,t,b} - \theta_{b_\ell,s,t,b}\right)
\end{equation}
where $X_\ell$ is the line reactance [$\Omega$] and $V_\ell$ is the
nominal voltage [kV].  A scaling factor $\sigma_\theta$
(\texttt{scale\_theta}, default 1000) is applied to the angle variables
for numerical conditioning: $\theta^\text{scaled} = \sigma_\theta \cdot
\theta$.

Line flow is bounded by transfer capacities:
\begin{equation}
-\overline{\phi}^{ba}_\ell \leq \phi_{\ell,s,t,b}
  \leq \overline{\phi}^{ab}_\ell
\end{equation}

\paragraph{Transmission losses.}
When line losses are enabled ($\lambda_\ell > 0$), separate forward and
reverse flow variables $\phi^+_\ell, \phi^-_\ell \geq 0$ replace the
single bidirectional variable.  The receiving bus sees only
$(1-\lambda_\ell)$ of the dispatched power, so losses are allocated at
the receiving end.  The Kirchhoff constraint~\eqref{eq:kirchhoff} then
uses $\phi^+_\ell - \phi^-_\ell$ in place of the single~$\phi_\ell$.

\paragraph{Transformer model.}
For tap-changing and phase-shifting transformers, the Kirchhoff
constraint generalizes to:
\begin{equation}
\label{eq:transformer}
-\theta_{a_\ell} + \theta_{b_\ell}
  + (\tau_\ell \chi_\ell)(\phi^+_\ell - \phi^-_\ell)
  = -\sigma_\theta \, \varphi_{\ell,\text{rad}}
\end{equation}
where $\tau_\ell$ is the off-nominal tap ratio (default~1),
$\chi_\ell = \sigma_\theta X_\ell / V_\ell^2$ is the scaled reactance,
and $\varphi_{\ell,\text{rad}}$ is the phase-shift angle in radians.
When $\tau_\ell = 1$ and $\varphi_\ell = 0$,~\eqref{eq:transformer}
reduces to~\eqref{eq:kirchhoff}.

\subsection{Battery Energy Storage}

State-of-charge balance for battery $e$ between consecutive blocks:
\begin{equation}
\label{eq:soc_balance}
e_{e,s,t,b} = e_{e,s,t,b-1}
  + p^\text{in}_{e,s,t,b} \cdot \eta^\text{in}_e \cdot \Delta_b
  - p^\text{out}_{e,s,t,b} \cdot \Delta_b / \eta^\text{out}_e
\end{equation}
where $\eta^\text{in}_e$ and $\eta^\text{out}_e$ are charge and discharge
efficiencies, and $p^\text{in}, p^\text{out}$ are charge and discharge
power [MW].  Initial SoC $e_{e,\cdot,\cdot,0} = e^\text{ini}_e$ is an
equality constraint in the first stage.

\begin{figure}[t]
  \centering
  \begin{tikzpicture}[
    >=Stealth,
    axis/.style={->, thick},
    charge/.style={fill=green!15},
    discharge/.style={fill=red!10},
  ]
    % Axes
    \draw[axis] (0,0) -- (7.2,0)
      node[right, font=\scriptsize] {Block (hour)};
    \draw[axis] (0,0) -- (0,3.2)
      node[above, font=\scriptsize] {SoC (\%)};

    % Y-axis ticks
    \foreach \y/\lbl in {0/0, 0.6/20, 1.2/40, 1.8/60, 2.4/80, 3.0/100} {
      \draw (-0.08,\y) -- (0.08,\y)
        node[left=2pt, font=\tiny] {\lbl};
    }
    % X-axis ticks
    \foreach \x/\lbl in
      {0.3/1, 1.5/6, 2.7/10, 3.6/13, 4.5/16, 5.7/20, 6.9/24} {
      \draw (\x,-0.08) -- (\x,0.08)
        node[below=2pt, font=\tiny] {\lbl};
    }

    % SoC data points
    \coordinate (p0) at (0.0, 0.9);   % hour 0:  30%
    \coordinate (p1) at (0.9, 0.6);   % hour 3:  20%
    \coordinate (p2) at (1.8, 0.6);   % hour 6:  20%
    \coordinate (p3) at (2.7, 1.5);   % hour 10: 50%
    \coordinate (p4) at (3.6, 2.7);   % hour 13: 90%
    \coordinate (p5) at (4.2, 2.7);   % hour 15: 90%
    \coordinate (p6) at (5.1, 1.5);   % hour 18: 50%
    \coordinate (p7) at (6.0, 0.9);   % hour 21: 30%
    \coordinate (p8) at (6.9, 0.9);   % hour 24: 30%

    % Charge shading (hours 6--15)
    \fill[charge] (p2) -- (p3) -- (p4) -- (p5)
      -- (p5 |- 0,0) -- (p2 |- 0,0) -- cycle;

    % Discharge shading (hours 15--21)
    \fill[discharge] (p5) -- (p6) -- (p7)
      -- (p7 |- 0,0) -- (p5 |- 0,0) -- cycle;

    % SoC line
    \draw[thick, blue!80] (p0) -- (p1) -- (p2) -- (p3) -- (p4)
      -- (p5) -- (p6) -- (p7) -- (p8);

    % Labels
    \node[font=\tiny, green!50!black, above] at (3.15,2.1)
      {charge};
    \node[font=\tiny, red!60!black, above] at (4.95,1.8)
      {discharge};

    % Solar peak annotation
    \draw[{Stealth[length=3pt]}-, thin, gray]
      (3.6,2.7) -- (4.3,3.1)
      node[right, font=\tiny] {solar peak};

    % Evening peak annotation
    \draw[{Stealth[length=3pt]}-, thin, gray]
      (5.1,1.5) -- (5.8,2.2)
      node[right, font=\tiny] {evening peak};
  \end{tikzpicture}
  \caption{Battery SoC over 24~blocks: charge during solar hours
    (green), discharge during evening peak (red).}
  \label{fig:battery_soc}
\end{figure}

\subsection{Converter Constraints}

A converter~$v$ couples a battery~$e$ to the electrical network through
a generator~$g$ (discharge path) and a demand~$d$ (charge path) via a
power conversion rate~$\rho_v$:
\begin{align}
\label{eq:conv_discharge}
p_{g,s,t,b} &= \rho_v \cdot p^\text{out}_{e,s,t,b} \\
\label{eq:conv_charge}
\ell_{d,s,t,b} &= \rho_v \cdot p^\text{in}_{e,s,t,b}
\end{align}
When the converter has a capacity limit~$\overline{C}_v$, a joint bound
applies: $p_{g,s,t,b} + \ell_{d,s,t,b} \leq \overline{C}_{v,t}$.
Converter capacity is expandable using the same modular framework
as~\eqref{eq:capacity}.

\subsection{Hydro Cascade}

For each junction $j$ with reservoir $r$, the water balance is:
\begin{equation}
\label{eq:hydro_balance}
v_{r,s,t,b} = v_{r,s,t,b-1}
  + \alpha_r \Delta_b \!\left(
    \textstyle\sum_{w \in \text{in}(j)} q_{w}
    - \textstyle\sum_{w \in \text{out}(j)} q_{w}
    + F_j
  \right)
\end{equation}
where $q_w$ is the waterway discharge [m$^3$/s], $F_j$ is the exogenous
inflow at junction $j$, and $\alpha_r = 0.0036$ converts flow rates to
volume increments: since 1~m$^3$/s flowing for 1~hour equals
3600~m$^3$ = 0.0036~hm$^3$, we have $\alpha_r = 3.6 \times 10^{-3}$
[hm$^3$/(m$^3$/s$\cdot$h)].

Each turbine $h$ linking waterway $w$ to generator $g$ provides:
\begin{equation}
\label{eq:turbine}
p_{g,s,t,b} \leq \rho_h \cdot q_{w,s,t,b}
\end{equation}
where $\rho_h$ is the water-to-power conversion rate [MW$\cdot$s/m$^3$].

\paragraph{Volume-dependent turbine efficiency.}
In practice, the conversion rate depends on the hydraulic head, which
varies with reservoir volume.  When defined, $\rho_h$ is replaced by a
concave piecewise-linear function:
\begin{equation}
\label{eq:vol_efficiency}
\rho_h(v) = \min_{i=1}^{N}
  \bigl\{ c_i + m_i (v - v_i) \bigr\}
\end{equation}
where $(v_i, m_i, c_i)$ are volume breakpoints, slopes, and intercepts
of~$N$ segments.  In the SDDP method, the active segment is updated at
each forward-pass iteration using the current reservoir volume.

\paragraph{Filtration (seepage).}
Seepage from a waterway into a reservoir is modeled as a linear function
of the average volume:
\begin{equation}
\label{eq:filtration}
\varphi_{i,s,t,b} = a_i + b_i \cdot
  \frac{v_{r,s,t,b-1} + v_{r,s,t,b}}{2}
\end{equation}
where $a_i$ is a constant term and $b_i$ is the seepage rate per unit
volume.  As with turbine efficiency, this admits a piecewise-linear
extension selected by the last known volume.

\begin{figure}[t]
  \centering
  \begin{tikzpicture}[
    res/.style={trapezium, trapezium angle=75, draw, fill=cyan!15,
      minimum width=1.4cm, minimum height=0.7cm, font=\scriptsize,
      trapezium stretches body},
    junc/.style={circle, draw, fill=gray!20, minimum size=0.4cm,
      font=\tiny, inner sep=1pt},
    turb/.style={regular polygon, regular polygon sides=3, draw,
      fill=yellow!25, minimum size=0.5cm, font=\tiny, inner sep=0pt},
    gen/.style={circle, draw, fill=orange!20, minimum size=0.5cm,
      font=\tiny},
    flow/.style={-{Stealth[length=4pt]}, thick, blue!70},
    spill/.style={-{Stealth[length=4pt]}, thick, red!50, dashed},
  ]
    % Reservoir 1
    \node[res] (r1) at (0,2.2) {$R_1$};
    \node[font=\tiny\itshape, text=black!70, right=1pt of r1] {$v_1$};

    % Junction 1
    \node[junc] (j1) at (0,1.1) {$j_1$};

    % Inflow
    \draw[flow] (-1.2,2.8) -- (j1.north west)
      node[midway, left, font=\tiny] {inflow $F_1$};

    % R1 -> J1
    \draw[flow] (r1) -- (j1);

    % Turbine
    \node[turb] (t1) at (1.2,0.6) {};
    \node[font=\tiny, right=1pt of t1] {$h_1$};

    % Junction 2
    \node[junc] (j2) at (0,-0.2) {$j_2$};

    % J1 -> T1
    \draw[flow] (j1) -- (t1) node[midway, above, font=\tiny] {$q_w$};

    % T1 -> Generator
    \node[gen] (g1) at (2.5,0.6) {\raisebox{-0.5pt}{\scalebox{0.5}{$\sim$}}};
    \node[font=\tiny, right=1pt of g1] {gen};
    \draw[-{Stealth[length=4pt]}, thick, orange!70] (t1) -- (g1);

    % J1 -> J2 spillway
    \draw[spill] (j1) -- (j2) node[midway, left, font=\tiny] {spill};

    % Reservoir 2
    \node[res] (r2) at (2,-0.2) {$R_2$};
    \node[font=\tiny\itshape, text=black!70, right=1pt of r2] {$v_2$};

    % Junction 3
    \node[junc] (j3) at (1,-1.2) {$j_3$};

    % J2 -> R2 waterway
    \draw[flow] (j2) -- (r2);

    % T1 -> J3 (tailwater)
    \draw[flow] (t1) |- (j3);

    % R2 -> J3
    \draw[flow] (r2) -- (j3);

    % Outflow
    \draw[flow] (j3) -- (1,-2.0) node[below, font=\tiny] {outflow};
  \end{tikzpicture}
  \caption{Hydro cascade example: two reservoirs ($R_1$, $R_2$),
    three junctions ($j_1$--$j_3$), one turbine ($h_1$) driving a
    generator, with inflow, spillway, and waterway connections.}
  \label{fig:hydro_cascade}
\end{figure}

\subsection{Capacity Expansion}

Installed capacity evolves as:
\begin{equation}
\label{eq:capacity}
C_{g,t} = C_{g,t-1}(1-\xi_g) + M_g \cdot m_{g,t}
\end{equation}
where $C_{g,0}$ is the initial installed capacity, $\xi_g$ is the annual
derating factor (\texttt{annual\_derating}), $M_g$ is the capacity per
expansion module (\texttt{expcap}), and $m_{g,t}$ is the number of modules
built.  The generator output is bounded by the installed capacity:
\begin{equation}
p_{g,s,t,b} \leq C_{g,t}
\end{equation}
Analogous expansion constraints apply to demands, lines, batteries, and
converters.

\subsection{Generator Profile Constraints}

When a generator has a capacity profile $\phi_g(s,t,b)$ (e.g., solar
irradiance or wind availability), the output is forced to follow the
profile:
\begin{equation}
\label{eq:profile}
p_{g,s,t,b} + s_{g,s,t,b} = \phi_g(s,t,b) \cdot C_{g,t}
\end{equation}
where $\phi_g(s,t,b) \in [0,1]$ is the per-block capacity factor and
$s_{g,s,t,b} \geq 0$ is the curtailed (spilled) generation.  The
spillover variable carries cost $c_g^\text{spill} \cdot \omega_{s,t,b}$,
which may be set to zero to allow free curtailment of renewables.  This
lets the optimizer spill excess renewable production when it is not
economically efficient to deliver.

\subsection{Demand Minimum Energy}

An optional minimum energy constraint ensures that a demand~$d$ is
served a minimum total energy within each stage:
\begin{equation}
\label{eq:min_energy}
\sum_{b \in \setB_t} \ell^\text{man}_{d,s,t,b} \cdot \Delta_b
  \geq E_d^\text{min}
\end{equation}
where $\ell^\text{man}_{d,s,t,b}$ is the mandatory served load and
$E_d^\text{min}$ [MWh] is the stage energy requirement.

\subsection{Reserve Constraints}

For each reserve zone $z$ and $(s,t,b)$:
\begin{equation}
\label{eq:reserve_up}
\sum_{g \in \mathcal{P}_z} u^\text{up}_{g,s,t,b}
  + u^\text{fail,up}_{z,s,t,b}
  \geq R^\text{up}_{z,s,t,b}
\end{equation}
where $u^\text{up}_{g,s,t,b}$ is the up-reserve provision from generator
$g$, $u^\text{fail,up}_{z,s,t,b}$ is unserved reserve, and
$R^\text{up}_{z,s,t,b}$ is the zone requirement.  A symmetric constraint
exists for down-reserve.  Reserve provision is bounded by the generator's
headroom: $u^\text{up}_{g,s,t,b} \leq C_{g,t} - p_{g,s,t,b}$.

\subsection{Constraint Summary}

Table~\ref{tab:constraints} lists all constraint families in the
formulation.

\begin{table}[t]
\centering
\caption{Complete constraint summary.}
\label{tab:constraints}
\small
\begin{tabular}{@{}clll@{}}
\toprule
\textbf{\#} & \textbf{Constraint} & \textbf{Eq.} & $\boldsymbol{\forall}$ \\
\midrule
C1  & Bus power balance             & \eqref{eq:balance}       & $n,s,t,b$ \\
C2  & Kirchhoff voltage law         & \eqref{eq:kirchhoff}     & $\ell,s,t,b$ \\
C3  & Generator output bounds       & \eqref{eq:pgen}          & $g,s,t,b$ \\
C4  & Capacity expansion            & \eqref{eq:capacity}      & $g,t$ \\
C5  & Generator profile             & \eqref{eq:profile}       & $g,s,t,b$ \\
C6  & Demand min.\ energy           & \eqref{eq:min_energy}    & $d,s,t$ \\
C7  & Line capacity                 & ---                      & $\ell,s,t,b$ \\
C8  & Battery SoC                   & \eqref{eq:soc_balance}   & $e,s,t,b$ \\
C9  & Converter coupling            & \eqref{eq:conv_discharge}& $v,s,t,b$ \\
C10 & Reserve requirements          & \eqref{eq:reserve_up}    & $z,s,t,b$ \\
C11 & Reserve--generator coupling   & ---                      & $p,s,t,b$ \\
C12 & Hydro water balance           & \eqref{eq:hydro_balance} & $j,s,t,b$ \\
C13 & Turbine conversion            & \eqref{eq:turbine}       & $h,s,t,b$ \\
C14 & Transformer / PST             & \eqref{eq:transformer}   & $\ell,s,t,b$ \\
\bottomrule
\end{tabular}
\end{table}
